%% =====================================================================
%% Journal manuscript (Elsevier elsarticle format), suitable for
%% submission to a Q2 ML/HEP journal (e.g. Nuclear Instruments and
%% Methods in Physics Research A, Computer Physics Communications,
%% Machine Learning: Science and Technology).
%% Build:  pdflatex -> bibtex -> pdflatex -> pdflatex
%% =====================================================================
\documentclass[preprint,12pt]{elsarticle}

\usepackage{amsmath,amssymb}
\usepackage{graphicx}
\usepackage{booktabs}
\usepackage{url}
\usepackage[hidelinks]{hyperref}

\graphicspath{{../../results/plots/}}

% Safe include: fall back to a placeholder if a figure is not present.
\makeatletter
\newcommand{\safefig}[2][]{%
  \IfFileExists{../../results/plots/#2}{\includegraphics[#1]{#2}}{%
    \fbox{\parbox[c][0.32\linewidth][c]{0.9\linewidth}{\centering
      \texttt{#2}\\[4pt]\textit{(figure generated by the pipeline)}}}}}
\makeatother

\journal{Computer Physics Communications}

\begin{document}

\begin{frontmatter}

%% --- Elsevier-style highlights (Q1 journals require 3--5, <=85 chars) ---
\begin{highlights}
\item Four top-taggers (BDT, CNN, ParticleNet, Particle Transformer) compared under one pipeline.
\item Data-efficiency sweep reveals a BDT-to-transformer crossover as training data grows.
\item Ranking and calibration are decoupled: the CNN is worst-calibrated despite mid ranking.
\item Temperature scaling fixes all models except the structurally miscalibrated CNN.
\item Fully reproducible on a single free cloud GPU; code and models released openly.
\end{highlights}

\title{Data efficiency and probability calibration of deep
jet-tagging architectures on the 14\,TeV top-quark benchmark}

\author[bnbwu]{Bibi Naifa}
\author[bnbwu]{Mah Lacka}
\author[bnbwu]{Kanwal Gul}
\author[bnbwu]{Firdos}
\author[bnbwu]{Ram Chand\corref{cor1}}
\ead{ram.chand2k11@yahoo.com}
\cortext[cor1]{Corresponding author.}

\affiliation[bnbwu]{organization={Department of Natural Sciences, The Begum
Nusrat Bhutto Women University},
            city={Sukkur},
            country={Pakistan}}

\begin{abstract}
Top tagging --- identifying the hadronic decays of highly boosted top
quarks --- is a standard benchmark for machine learning in collider
physics. We compare four classifiers on the public Top-Quark Tagging
Reference Dataset of 14\,TeV proton--proton collisions: a gradient-boosted
decision tree (BDT) on engineered jet-substructure variables, a
convolutional neural network (CNN) on jet images, the ParticleNet graph
network, and the Particle Transformer. Rather than chasing a new record
area under the curve (AUC), we study two properties of direct relevance to
downstream analyses: \emph{data efficiency} and \emph{probability
calibration}. Using a fully reproducible pipeline trained on a single free
cloud GPU, we sweep the training-set size from $5\times10^{3}$ to
$5\times10^{4}$ jets and measure test AUC, background rejection at fixed
signal efficiency, the Expected Calibration Error (ECE), and the effect of
temperature scaling. The Particle Transformer attains the best
discrimination (AUC $=0.978$), reproducing the published ordering, but the
data-efficiency sweep reveals a clear crossover: the BDT is strongest at the
smallest training size and is only overtaken by the deep models as data
grows. The calibration study shows that ranking power and calibration are
decoupled --- the jet-image CNN is the worst-calibrated model (ECE $=0.078$)
despite a mid-pack ranking, whereas the top-ranked transformer is
comparatively well-calibrated, and a single-parameter temperature rescaling
corrects every model except the CNN. The complete pipeline, trained models
and documentation are released as open source.
\end{abstract}

\begin{keyword}
top tagging \sep jet substructure \sep deep learning \sep ParticleNet \sep
Particle Transformer \sep data efficiency \sep probability calibration
\end{keyword}

\end{frontmatter}

%% =====================================================================
\section{Introduction}
\label{sec:intro}

\subsection{Physics context}
The top quark is the heaviest known elementary particle, with a mass
$m_t\simeq172.76$\,GeV \citep{pdg2022}. Because this mass exceeds the
hadronisation scale, the top decays as a bare quark, almost always through
$t\rightarrow Wb$; in the fully hadronic channel the $W$ decays as
$W\rightarrow q\bar{q}'$, giving $t\rightarrow Wb\rightarrow q\bar{q}'b$. When
the parent top carries a transverse momentum $p_T\gg m_t$ it is
\emph{boosted}: its three decay quarks are produced collinearly and, after
parton showering and hadronisation, merge into one large-radius jet
reconstructed by sequential-recombination algorithms such as anti-$k_t$
\citep{cacciari2008,cacciari2012fastjet,salam2010}. Such a ``top jet'' differs
from an ordinary jet initiated by a light quark or gluon (a QCD jet) in three
measurable ways: it is more massive (mass $\sim m_t$), it shows a
three-pronged rather than one-pronged radiation pattern, and it carries a
higher constituent multiplicity. Exploiting these differences to separate top
jets from the overwhelmingly more abundant QCD background --- \emph{top
tagging} --- is one of the canonical classification problems of LHC physics
and a central tool in searches for physics beyond the Standard Model
\citep{larkoski2020,kogler2019}.

\subsection{From engineered features to deep learning}
Classical taggers compress the jet into a handful of physically motivated
\emph{substructure observables}: the jet mass, the $N$-subjettiness ratios
$\tau_{21},\tau_{32}$ \citep{thaler2011}, and energy-correlation functions
such as $C_2$ and $D_2$ \citep{larkoski2013,larkoski2014}. Fed to a
boosted decision tree (BDT), these provide a strong, interpretable baseline.
The deep-learning era replaced hand-crafted features with representations
learned directly from the constituents. Three paradigms followed in
succession: pixelating the jet into a calorimeter-like image for a
convolutional neural network (CNN) \citep{deoliveira2016,komiske2017};
treating the jet as an unordered ``particle cloud'' processed by a
permutation-invariant graph network such as ParticleNet
\citep{qu2020particlenet}, built on the Dynamic Graph CNN \citep{wang2019dgcnn};
and, most recently, modelling pairwise constituent interactions with
self-attention in the Particle Transformer
\citep{qu2022particletransformer,vaswani2017attention}. Related set, graph and
Lorentz-equivariant architectures continue to advance the field
\citep{komiske2019efn,moreno2020jedinet,bogatskiy2020lorentz,gong2022lorentznet}.
On the standard reference dataset \citep{kasieczka2019dataset}, a community
comparison \citep{kasieczka2019landscape} established the now-canonical
ordering, with the leading deep models exceeding an area under the receiver
operating characteristic curve (AUC) of $0.98$.

\subsection{Motivation and contributions}
The published literature is almost entirely concerned with the \emph{best
achievable} discrimination, quoted at the full training-set size of
$1.2\times10^{6}$ jets and on large GPU clusters. Two properties of direct
practical importance are rarely studied together. The first is \emph{data
efficiency}: generating labelled Monte-Carlo training samples is
computationally expensive, so how steeply each architecture's performance
depends on the amount of training data --- and whether the ranking of the
models changes in the small-data regime --- is decisive for groups without
large simulation budgets. The second is \emph{probability calibration}: a high
AUC certifies only that a model \emph{ranks} signal above background, whereas
many downstream analyses use the classifier score as a \emph{probability}
(for instance in a profile-likelihood fit), which requires the score to be
calibrated. While calibration is routinely examined in the machine-learning
literature \citep{guo2017calibration,niculescu2005calibration}, it is seldom
reported for jet taggers.

This paper addresses both questions on the public reference dataset using a
single, reproducible pipeline that runs end-to-end on free hardware. Our
contributions are:
\begin{enumerate}
    \item a unified, open-source pipeline that trains and evaluates the BDT,
    CNN, ParticleNet and Particle Transformer under identical preprocessing,
    data splits and random seeds;
    \item a data-efficiency study across a decade of training-set sizes
    ($5\times10^{3}$ to $5\times10^{4}$ jets), revealing a clear crossover
    between the BDT and the deep models; and
    \item a probability-calibration study, using reliability diagrams, the
    Expected Calibration Error and temperature scaling, showing that ranking
    power and calibration are decoupled, with practical recommendations for
    downstream use.
\end{enumerate}
The remainder of the paper is organised as follows.
Section~\ref{sec:data} describes the dataset and preprocessing;
Section~\ref{sec:methodology} sets out the four models and the evaluation
methodology with their governing equations; Section~\ref{sec:results}
presents the headline benchmark, the data-efficiency sweep and the
calibration analysis; and Section~\ref{sec:conclusions} concludes.

%% =====================================================================
\section{Dataset and preprocessing}
\label{sec:data}

We use the public Top-Quark Tagging Reference Dataset
\citep{kasieczka2019dataset,kasieczka2019landscape}, $2\times10^{6}$ jets from
14\,TeV $pp$ collisions generated with \textsc{Pythia\,8}
\citep{sjostrand2015pythia} and a \textsc{Delphes} fast detector simulation
\citep{delphes2014}, clustered with the anti-$k_t$ algorithm
\citep{cacciari2008,cacciari2012fastjet} at $R=0.8$ and $550<p_T<650$\,GeV.
For tractability we draw stratified subsamples of up to $5\times10^{4}$ jets
per split and use the leading $100$ constituents per jet.

For each constituent we form the feature vector
$(p_T,\eta,\phi,E,\Delta\eta,\Delta\phi,\log p_T)$ relative to the
$p_T$-weighted jet axis, standardised on the training split. For the BDT we
compute thirteen engineered observables (jet mass, $p_T$, $\eta$,
multiplicity, width, leading/sub-leading $p_T$ fractions,
$\tau_1,\tau_2,\tau_3$, $\tau_{21},\tau_{32}$ \citep{thaler2011}, and $C_2$
\citep{larkoski2013,larkoski2014}); their distributions
(Fig.~\ref{fig:substructure}) behave exactly as physics expects --- the top
mass peaks at $m_t$ and $\tau_{32}$ separates the three-prong signal. For the
CNN we bin constituents into $40\times40\times3$ jet images
\citep{deoliveira2016}; the averaged images (Fig.~\ref{fig:avgimg}) show the
three-prong topology in the signal$-$background difference. For the Particle
Transformer we add four pairwise features
$(\ln\Delta R,\ln k_T,z,\Delta\eta)$ following
\citet{qu2022particletransformer}.

\begin{figure}[t]
\centering
\safefig[width=\linewidth]{jet_substructure.pdf}
\caption{Engineered jet-substructure observables for top (red) and QCD
(blue) jets on the test set.}
\label{fig:substructure}
\end{figure}

\begin{figure}[t]
\centering
\safefig[width=\linewidth]{average_jet_images.pdf}
\caption{Averaged $p_T$-weighted jet images: top (left), QCD (centre) and
their difference (right).}
\label{fig:avgimg}
\end{figure}

%% =====================================================================
\section{Methodology}
\label{sec:methodology}

We treat top tagging as binary classification: each jet $\mathbf{x}$ is mapped
by a model $f_\theta$ to a score $\hat{y}=f_\theta(\mathbf{x})\in[0,1]$
estimating the probability that the jet is a top jet. All four models are
trained against the binary cross-entropy loss
\begin{equation}
\mathcal{L}(\theta) = -\frac{1}{N}\sum_{i=1}^{N}
\Bigl[\, y_i\log\hat{y}_i + (1-y_i)\log(1-\hat{y}_i) \,\Bigr],
\label{eq:bce}
\end{equation}
where $y_i\in\{0,1\}$ is the truth label and the neural models produce
$\hat{y}_i=\sigma(z_i)$ from a logit $z_i$ through the logistic sigmoid
$\sigma(z)=1/(1+e^{-z})$.

\subsection{Boosted decision tree (BDT)}
\label{sec:bdt}
The baseline is a gradient-boosted decision-tree ensemble trained with
\textsc{XGBoost} \citep{chen2016xgboost} on the thirteen engineered
observables. The prediction is an additive ensemble of $K$ regression trees
$f_k\in\mathcal{F}$,
\begin{equation}
z_i = \sum_{k=1}^{K} f_k(\mathbf{x}_i),
\label{eq:bdt-ensemble}
\end{equation}
built greedily by minimising the regularised objective
\begin{equation}
\mathcal{O} = \sum_{i} \ell(y_i,\hat{y}_i)
            + \sum_{k}\Omega(f_k), \qquad
\Omega(f) = \gamma\,T + \tfrac{1}{2}\lambda\lVert\mathbf{w}\rVert^{2},
\label{eq:xgb-obj}
\end{equation}
where $\ell$ is the logistic loss, $T$ is the number of leaves, $\mathbf{w}$
the leaf weights, and $\gamma,\lambda$ control complexity. At each boosting
round a new tree is fit to the second-order Taylor expansion of $\mathcal{O}$,
which yields fast, accurate splits. Among the inputs, the $N$-subjettiness
ratios \citep{thaler2011}
\begin{equation}
\tau_N = \frac{1}{d_0}\sum_{k} p_{T,k}\,
\min\bigl(\Delta R_{1,k},\dots,\Delta R_{N,k}\bigr), \qquad
\tau_{21}=\frac{\tau_2}{\tau_1},\;\; \tau_{32}=\frac{\tau_3}{\tau_2},
\label{eq:nsub}
\end{equation}
($d_0=\sum_k p_{T,k}R_0$, $\Delta R_{a,k}$ the angular distance from
constituent $k$ to axis $a$) measure compatibility with an $N$-pronged
hypothesis and are, with the jet mass, the most discriminating features.

\subsection{Convolutional neural network (CNN)}
\label{sec:cnn}
The CNN acts on the $40\times40\times3$ jet image $I$, applying learnable
kernels $W$ in the $(\Delta\eta,\Delta\phi)$ plane,
\begin{equation}
(I * W)_{ij}^{c'} = \sum_{c}\sum_{a,b} W^{c',c}_{a,b}\,
I^{c}_{\,i+a,\,j+b},
\label{eq:conv}
\end{equation}
interleaved with batch normalisation and ReLU nonlinearities
$\mathrm{ReLU}(u)=\max(0,u)$. We use a ResNet-style stack
\citep{he2016resnet} in which each block learns a residual mapping
\begin{equation}
\mathbf{h}_{\ell+1} = \mathrm{ReLU}\bigl(\mathbf{h}_\ell
+ \mathcal{F}(\mathbf{h}_\ell;W_\ell)\bigr),
\label{eq:resnet}
\end{equation}
easing optimisation of deeper networks. Global average pooling and a
fully connected head produce the logit $z$.

\subsection{ParticleNet (graph network)}
\label{sec:particlenet}
ParticleNet \citep{qu2020particlenet} treats the jet as an unordered point
cloud $\{\mathbf{x}_i\}$ and is built from EdgeConv blocks
\citep{wang2019dgcnn}. For each particle $i$, a $k$-nearest-neighbour graph is
constructed in coordinate space and the feature is updated by a symmetric
aggregation over its neighbourhood $\mathcal{N}(i)$,
\begin{equation}
\mathbf{x}_i' = \operatorname*{max}_{j\in\mathcal{N}(i)}\;
h_\Theta\!\bigl(\mathbf{x}_i,\;\mathbf{x}_j-\mathbf{x}_i\bigr),
\label{eq:edgeconv}
\end{equation}
where $h_\Theta$ is a shared multilayer perceptron and the element-wise
$\max$ makes the operation permutation-invariant. The graph is recomputed in
the learned feature space at each block (a ``dynamic'' graph), so successive
layers capture progressively larger-scale structure; a final global pooling
and dense head give $z$.

\subsection{Particle Transformer}
\label{sec:part}
The Particle Transformer \citep{qu2022particletransformer} applies multi-head
self-attention \citep{vaswani2017attention} to the constituent sequence. With
queries, keys and values $Q,K,V$ obtained by linear projection, a standard
attention head computes
\begin{equation}
\mathrm{Attention}(Q,K,V) =
\mathrm{softmax}\!\left(\frac{QK^{\top}}{\sqrt{d_k}}\right)V,
\label{eq:attention}
\end{equation}
with $d_k$ the key dimension. The key physics ingredient is a learned
\emph{interaction bias}: the four pairwise features
$\mathbf{u}_{ij}=(\ln\Delta R_{ij},\ln k_{T,ij},z_{ij},\Delta\eta_{ij})$ are
mapped by an MLP to a per-head additive term $U_{ij}=\phi(\mathbf{u}_{ij})$
that is injected into the attention logits,
\begin{equation}
A = \mathrm{softmax}\!\left(\frac{QK^{\top}}{\sqrt{d_k}} + U\right)V,
\label{eq:part-attention}
\end{equation}
so that physically close, hard splittings are emphasised. A learnable class
token aggregates the sequence and its representation is mapped to the logit
$z$; its attention weights provide an interpretability handle
(Section~\ref{sec:results}).

\subsection{Training}
\label{sec:training}
The neural models minimise Eq.~\eqref{eq:bce} with Adam/AdamW
\citep{kingma2014adam,loshchilov2019adamw} (learning rates $10^{-3}$--$10^{-4}$),
gradient clipping, a plateau learning-rate scheduler and early stopping on the
validation AUC, implemented in \textsc{PyTorch} \citep{paszke2019pytorch} with
\textsc{NumPy} \citep{harris2020numpy}, \textsc{scikit-learn}
\citep{pedregosa2011scikit} and \textsc{Matplotlib}
\citep{hunter2007matplotlib}. Training used a free Google Colaboratory T4 GPU;
the pipeline checkpoints every epoch and resumes automatically.

\subsection{Evaluation metrics}
\label{sec:metrics}
We report the AUC, the accuracy at a score threshold of $0.5$, and the
background rejection $1/\varepsilon_B$ at fixed signal efficiencies
$\varepsilon_S=0.5,0.3$, the figure of merit most relevant to physics
analyses; AUC uncertainties are obtained by bootstrap resampling of the test
set. Calibration is quantified by the Expected Calibration Error
\citep{guo2017calibration}, which partitions the predictions into $M$
equal-width score bins $B_m$ and averages the gap between confidence and
accuracy,
\begin{equation}
\mathrm{ECE} = \sum_{m=1}^{M}\frac{|B_m|}{N}\,
\bigl|\,\mathrm{acc}(B_m)-\mathrm{conf}(B_m)\,\bigr|.
\label{eq:ece}
\end{equation}
As a corrective we apply \emph{temperature scaling}: a single scalar $T>0$ is
fit on the validation set by minimising the negative log-likelihood of the
rescaled scores $\hat{y}=\sigma(z/T)$, and we report the ECE before and after
rescaling. $T>1$ indicates an over-confident model whose probabilities are
softened by the rescaling.

%% =====================================================================
\section{Results}
\label{sec:results}

\subsection{Headline benchmark}
Table~\ref{tab:comparison} and Fig.~\ref{fig:roc} report the headline
performance. The ordering Particle Transformer ($0.978$) $>$ ParticleNet
($0.976$) $>$ CNN ($0.972$) $>$ BDT ($0.970$) reproduces the community
benchmark \citep{kasieczka2019landscape}; absolute values are slightly below
the published saturation values because we train on a subset rather than the
full $1.2\times10^{6}$ jets.

\begin{table}[t]
\centering
\caption{Headline test-set performance ($5\times10^{4}$-jet training run).}
\label{tab:comparison}
\small
\begin{tabular}{lcccc}
\toprule
Model & AUC & Acc. & $1/\varepsilon_B^{0.5}$ & $1/\varepsilon_B^{0.3}$ \\
\midrule
BDT (XGBoost)        & 0.970 & 0.911 & 93  & 287 \\
CNN (jet images)     & 0.972 & 0.915 & 111 & 413 \\
ParticleNet          & 0.976 & 0.920 & 130 & 413 \\
Particle Transformer & 0.978 & 0.925 & 136 & 476 \\
\bottomrule
\end{tabular}
\end{table}

\begin{figure}[t]
\centering
\safefig[width=\linewidth]{roc_comparison.pdf}
\caption{ROC curves (left) and background rejection vs.\ signal efficiency
(right).}
\label{fig:roc}
\end{figure}

\subsection{Data efficiency}
Re-training each model at
$N_{\text{train}}\in\{5,10,25,50\}\times10^{3}$ jets yields the learning
curves of Fig.~\ref{fig:learning}. The BDT is the most data-efficient (AUC
$0.961\rightarrow0.968$ across the range), whereas the deep models gain
strongly and \emph{cross over}: the Particle Transformer is the weakest model
at $5\times10^{3}$ jets (AUC $0.949$) but the strongest at $5\times10^{4}$
($0.974$). Under a tight Monte-Carlo budget the simple, interpretable BDT is
therefore competitive or preferable; the attention model only realises its
advantage with sufficient data, consistent with its near-$10^{6}$-jet
saturation \citep{qu2022particletransformer}.

\begin{figure}[t]
\centering
\safefig[width=\linewidth]{learning_curves.pdf}
\caption{Data-efficiency curves: test AUC (left) and background rejection at
$\varepsilon_S=0.5$ (right) vs.\ training-set size.}
\label{fig:learning}
\end{figure}

\begin{table}[t]
\centering
\caption{Test AUC versus training-set size $N_{\text{train}}$. Uncertainties
are one standard deviation from a $200$-sample bootstrap of the test set. The
Particle Transformer is weakest at $5\times10^{3}$ jets and strongest at
$5\times10^{4}$, crossing the BDT and ParticleNet.}
\label{tab:efficiency}
\small
\setlength{\tabcolsep}{5pt}
\begin{tabular}{lcccc}
\toprule
Model & $5\!\times\!10^{3}$ & $10^{4}$ & $2.5\!\times\!10^{4}$ & $5\!\times\!10^{4}$ \\
\midrule
BDT (XGBoost)        & 0.961$\pm$0.001 & 0.964$\pm$0.001 & 0.967$\pm$0.001 & 0.968$\pm$0.001 \\
CNN (jet images)     & 0.947$\pm$0.001 & 0.953$\pm$0.001 & 0.968$\pm$0.001 & 0.968$\pm$0.001 \\
ParticleNet          & 0.950$\pm$0.001 & 0.955$\pm$0.001 & 0.966$\pm$0.001 & 0.971$\pm$0.001 \\
Particle Transformer & 0.949$\pm$0.001 & 0.952$\pm$0.001 & 0.966$\pm$0.001 & 0.974$\pm$0.001 \\
\bottomrule
\end{tabular}
\end{table}

\subsection{Probability calibration}
Following \citet{guo2017calibration} we compute the Expected Calibration
Error (ECE) over $15$ score bins and fit a single temperature $T$ on the
validation set. Figure~\ref{fig:calib} and Table~\ref{tab:calib} show that
all models are mildly over-confident ($T>1$) but to very different degrees.
Strikingly, ranking and calibration are \emph{decoupled}: the jet-image CNN
is the worst-calibrated by an order of magnitude (ECE $=0.078$) despite
ranking third on AUC, whereas the top-ranked Particle Transformer is
comparatively well-calibrated (ECE $=0.012$). Temperature scaling brings the
BDT, ParticleNet and transformer to ECE $\lesssim0.008$ but only halves the
CNN's error (to $0.064$): the CNN's miscalibration is structural and would
require a richer recalibration such as isotonic regression
\citep{niculescu2005calibration}.

\begin{figure}[t]
\centering
\safefig[width=\linewidth]{calibration.pdf}
\caption{Reliability diagrams (top) and score histograms (bottom). Solid:
raw scores; dotted: temperature-scaled. The CNN deviates most.}
\label{fig:calib}
\end{figure}

\begin{table}[t]
\centering
\caption{Calibration at the $5\times10^{4}$-jet endpoint: ECE, fitted
temperature $T$, and ECE after rescaling.}
\label{tab:calib}
\small
\begin{tabular}{lccc}
\toprule
Model & ECE & $T$ & ECE$_T$ \\
\midrule
BDT (XGBoost)        & 0.007 & 1.04 & 0.006 \\
CNN (jet images)     & 0.078 & 1.54 & 0.064 \\
ParticleNet          & 0.008 & 1.05 & 0.008 \\
Particle Transformer & 0.012 & 1.12 & 0.007 \\
\bottomrule
\end{tabular}
\end{table}

\subsection{Interpretability}
The class-token attention of the Particle Transformer
(Fig.~\ref{fig:clsattn}) concentrates at small angular distance $\Delta R$
for QCD jets (a single core) but spreads to large $\Delta R$ for top jets,
reflecting the broader multi-prong energy flow of the top decay. The BDT
feature importances are dominated by the jet mass, $\tau_{32}$ and
multiplicity, the physically expected discriminators.

\begin{figure}[t]
\centering
\safefig[width=\linewidth]{cls_attention.pdf}
\caption{Class-token attention vs.\ constituent $\Delta R$ for top jets
(left) and QCD jets (right).}
\label{fig:clsattn}
\end{figure}

%% =====================================================================
\section{Conclusions}
\label{sec:conclusions}

We have presented a reproducible comparison of four top-tagging
architectures focused on data efficiency and probability calibration. The
literature ordering is reproduced, but two practical messages emerge. First,
data efficiency separates the models: the BDT is best in the small-data
regime, with the deep models overtaking it only as data grows. Second,
ranking and calibration are decoupled --- the best-ranked transformer is
well-calibrated, while the CNN, mid-pack on AUC, is the worst-calibrated and
is not fully corrected by temperature scaling. We recommend reporting
temperature-scaled probabilities by default and treating CNN scores with
particular caution. Natural extensions are to scale the sweep to the full
dataset, to study robustness under pile-up, and to compare non-parametric
recalibration.

%% =====================================================================
\section*{Code and data availability}
The complete pipeline, trained models, and documents are openly available at
\url{https://github.com/ramc77/fyp-jet-tagging}. The dataset is the public
Top-Quark Tagging Reference Dataset \citep{kasieczka2019dataset}
(\url{https://zenodo.org/records/2603256}).

\section*{CRediT authorship contribution statement}
\textbf{Bibi Naifa:} Methodology, Software, Investigation, Writing -- original
draft. \textbf{Mah Lacka:} Software, Validation, Visualization, Writing --
original draft. \textbf{Kanwal Gul:} Formal analysis, Data curation, Writing
-- review \& editing. \textbf{Firdos:} Investigation, Validation, Writing --
review \& editing. \textbf{Ram Chand:} Conceptualization, Supervision, Project
administration, Methodology, Writing -- review \& editing.

\section*{Declaration of competing interest}
The authors declare that they have no known competing financial interests or
personal relationships that could have appeared to influence the work
reported in this paper.

\section*{Acknowledgements}
We thank the maintainers of the Top-Quark Tagging Reference Dataset and the
open-source scientific-Python and Google Colaboratory ecosystems. This
research received no specific grant from any funding agency.

%% Ensure supporting software/method references appear in the list.
\nocite{salam2010,komiske2019efn,paszke2019pytorch,harris2020numpy,
pedregosa2011scikit,hunter2007matplotlib}

\bibliographystyle{elsarticle-num-names}
\bibliography{refs}

\end{document}
